\begin{corollary}[逼近 \texorpdfstring{$T\log T$}{T log T} 临界线密度的状态规模下界]\label{cor:critical-circle-state-dimension-lower-bound}
沿用命题 \ref{prop:critical-circle-density-calibration} 与推论 \ref{cor:critical-circle-mode-growth} 的记号。
设 \(M\) 为一个 \(\,N\times N\) 的有限维核矩阵，并令 \(\lambda=\rho(M)>1\)。
记临界圆模态数 \(B\) 为满足 \(|\mu|=\sqrt{\lambda}\) 的特征值（按代数重数计）个数，则必有
\[
\boxed{\;B\le N.\;}
\]
因此若希望在高度 \(T\) 内达到 \(N_{\mathrm{crit}}(T)\asymp T\log T\) 的临界线零极点密度，则任何固定 \(\lambda>1\) 的有限维核族 \(M_T\) 都必须满足
\[
\boxed{\;N(T)\ \ge\ B(T)\ \ge\ (1-o(1))\frac{\log T}{\log\lambda}\;}
\qquad(T\to\infty),
\]
从而任何固定有限状态数（常数维度）的核模型只能产生 \(O(T)\) 的线性密度，无法在计数级别匹配 Riemann--von Mangoldt 主项的 \(T\log T\) 规模。
\end{corollary}

\begin{proof}
由于 \(M\) 的全部特征值（按代数重数计）总数为 \(N\)，而临界圆模态 \(\{\mu_b\}_{b=1}^B\) 仅为其中一部分，故 \(B\le N\)。
其余结论由推论 \ref{cor:critical-circle-mode-growth} 的必要条件 \(B(T)\asymp (\log T)/(\log\lambda)\) 直接推出。证毕。
\end{proof}

\endinput
